% Source: book pp. 228--242 (PDF 242--256). No solutions in the manual.
% a numbered display (Axler numbers a few equations with the item counter):
%   text\axeqstep{ax:7.15}  \[ ... \tag{\thetheorem} \]
\DeclareRobustCommand\axeqstep[1]{\refstepcounter{theorem}\label{#1}}
% the small italic comment that follows some exercises
\section{Self-Adjoint and Normal Operators}

\subsection{Adjoints}

\begin{definition}[Adjoint, $T^*$]\label{ax:7.1}
Suppose $T\in\Lin(V,W)$. The \emph{adjoint} of $T$ is the function
$T^*\colon W\to V$ such that
\[ \ip{Tv}{w} = \ip{v}{T^*w} \]
for every $v\in V$ and every $w\in W$.
\end{definition}

\marginnote{The word \textbf{adjoint} has another meaning in linear algebra. In
case you encounter the second meaning elsewhere, be warned that the two meanings
for adjoint are unrelated to each other.}%
To see why the definition above makes sense, suppose $T\in\Lin(V,W)$. Fix
$w\in W$. Consider the linear functional
\[ v\mapsto\ip{Tv}{w} \]
on $V$ that maps $v\in V$ to $\ip{Tv}{w}$; this linear functional depends on $T$
and $w$. By the Riesz representation theorem (\ref{ax:6.42}), there exists a
unique vector in $V$ such that this linear functional is given by taking the
inner product with it. We call this unique vector $T^*w$. In other words, $T^*w$
is the unique vector in $V$ such that
\[ \ip{Tv}{w} = \ip{v}{T^*w} \]
for every $v\in V$.

In the equation above, the inner product on the left takes place in $W$ and the
inner product on the right takes place in $V$. However, we use the same
notation $\ip{\cdot}{\cdot}$ for both inner products.

\begin{example}[Adjoint of a linear map from $\R^3$ to $\R^2$]\label{ax:7.2}
Define $T\colon\R^3\to\R^2$ by
\[ T(x_1,x_2,x_3) = (x_2+3x_3,\,2x_1). \]
To compute $T^*$, suppose $(x_1,x_2,x_3)\in\R^3$ and $(y_1,y_2)\in\R^2$. Then
\begin{align*}
\bigl\langle T(x_1,x_2,x_3),(y_1,y_2)\bigr\rangle
  &= \bigl\langle (x_2+3x_3,\,2x_1),(y_1,y_2)\bigr\rangle\\
  &= x_2y_1+3x_3y_1+2x_1y_2\\
  &= \bigl\langle (x_1,x_2,x_3),(2y_2,y_1,3y_1)\bigr\rangle.
\end{align*}
The equation above and the definition of the adjoint imply that
\[ T^*(y_1,y_2) = (2y_2,y_1,3y_1). \]
\end{example}

\begin{example}[Adjoint of a linear map with range of dimension at most $1$]%
\label{ax:7.3}
Fix $u\in V$ and $x\in W$. Define $T\in\Lin(V,W)$ by
\[ Tv = \ip{v}{u}x \]
for each $v\in V$. To compute $T^*$, suppose $v\in V$ and $w\in W$. Then
\begin{align*}
\ip{Tv}{w} &= \bigl\langle \ip{v}{u}x,w\bigr\rangle\\
           &= \ip{v}{u}\ip{x}{w}\\
           &= \bigl\langle v,\ip{w}{x}u\bigr\rangle.
\end{align*}
Thus
\[ T^*w = \ip{w}{x}u. \]
\end{example}

\marginnote[-3em]{The two examples above and the proof below use a common technique
for computing $T^*$: start with a formula for $\ip{Tv}{w}$ then manipulate it to
get just $v$ in the first slot; the entry in the second slot will then be
$T^*w$.}%
In the two examples above, $T^*$ turned out to be not just a function from $W$
to $V$ but a linear map from $W$ to $V$. This behavior is true in general, as
shown by the next result.

\begin{result}[Adjoint of a linear map is a linear map]\label{ax:7.4}
If $T\in\Lin(V,W)$, then $T^*\in\Lin(W,V)$.
\end{result}

\begin{proof}
Suppose $T\in\Lin(V,W)$. If $v\in V$ and $w_1,w_2\in W$, then
\begin{align*}
\ip{Tv}{w_1+w_2} &= \ip{Tv}{w_1}+\ip{Tv}{w_2}\\
                 &= \ip{v}{T^*w_1}+\ip{v}{T^*w_2}\\
                 &= \ip{v}{T^*w_1+T^*w_2}.
\end{align*}
The equation above shows that
\[ T^*(w_1+w_2) = T^*w_1+T^*w_2. \]

If $v\in V$, $\lambda\in\F$, and $w\in W$, then
\begin{align*}
\ip{Tv}{\lambda w} &= \overline{\lambda}\ip{Tv}{w}\\
                   &= \overline{\lambda}\ip{v}{T^*w}\\
                   &= \ip{v}{\lambda T^*w}.
\end{align*}
The equation above shows that
\[ T^*(\lambda w) = \lambda T^*w. \]

Thus $T^*$ is a linear map, as desired.
\end{proof}

\begin{result}[Properties of the adjoint]\label{ax:7.5}
Suppose $T\in\Lin(V,W)$. Then
\begin{parts}
\item $(S+T)^*=S^*+T^*$ for all $S\in\Lin(V,W)$;
\item $(\lambda T)^*=\overline{\lambda}T^*$ for all $\lambda\in\F$;
\item $(T^*)^*=T$;
\item $(ST)^*=T^*S^*$ for all $S\in\Lin(W,U)$ (here $U$ is a finite-dimensional
  inner product space over $\F$);
\item $I^*=I$, where $I$ is the identity operator on $V$;
\item if $T$ is invertible, then $T^*$ is invertible and
  $(T^*)^{-1}=(T^{-1})^*$.
\end{parts}
\end{result}

\begin{proof}
Suppose $v\in V$ and $w\in W$.
\begin{parts}
\item If $S\in\Lin(V,W)$, then
  \begin{align*}
  \bigl\langle (S+T)v,w\bigr\rangle &= \ip{Sv}{w}+\ip{Tv}{w}\\
    &= \ip{v}{S^*w}+\ip{v}{T^*w}\\
    &= \ip{v}{S^*w+T^*w}.
  \end{align*}
  Thus $(S+T)^*w=S^*w+T^*w$, as desired.
\item If $\lambda\in\F$, then
  \[ \bigl\langle (\lambda T)v,w\bigr\rangle = \lambda\ip{Tv}{w}
     = \lambda\ip{v}{T^*w} = \ip{v}{\overline{\lambda}T^*w}. \]
  Thus $(\lambda T)^*w=\overline{\lambda}T^*w$, as desired.
\item We have
  \[ \ip{T^*w}{v} = \overline{\ip{v}{T^*w}} = \overline{\ip{Tv}{w}}
     = \ip{w}{Tv}. \]
  Thus $(T^*)^*v=Tv$, as desired.
\item Suppose $S\in\Lin(W,U)$ and $u\in U$. Then
  \[ \bigl\langle (ST)v,u\bigr\rangle = \bigl\langle S(Tv),u\bigr\rangle
     = \ip{Tv}{S^*u} = \bigl\langle v,T^*(S^*u)\bigr\rangle. \]
  Thus $(ST)^*u=T^*(S^*u)$, as desired.
\item Suppose $u\in V$. Then
  \[ \ip{Iu}{v} = \ip{u}{v}. \]
  Thus $I^*v=v$, as desired.
\item Suppose $T$ is invertible. Take adjoints of both sides of the equation
  $T^{-1}T=I$, then use (d) and (e) to show that $T^*(T^{-1})^*=I$. Similarly,
  the equation $TT^{-1}=I$ implies $(T^{-1})^*T^*=I$. Thus $(T^{-1})^*$ is the
  inverse of $T^*$, as desired. \qedhere
\end{parts}
\end{proof}

If $\F=\R$, then the map $T\mapsto T^*$ is a linear map from $\Lin(V,W)$ to
$\Lin(W,V)$, as follows from (a) and (b) of the result above. However, if
$\F=\C$, then this map is not linear because of the complex conjugate that
appears in (b).

The next result shows the relationship between the null space and the range of
a linear map and its adjoint.

\begin{result}[Null space and range of $T^*$]\label{ax:7.6}
Suppose $T\in\Lin(V,W)$. Then
\begin{parts}
\item $\nul T^*=(\range T)^\perp$;
\item $\range T^*=(\nul T)^\perp$;
\item $\nul T=(\range T^*)^\perp$;
\item $\range T=(\nul T^*)^\perp$.
\end{parts}
\end{result}

\begin{proof}
We begin by proving (a). Let $w\in W$. Then
\begin{align*}
w\in\nul T^* &\iff T^*w=0\\
  &\iff \ip{v}{T^*w}=0 \text{ for all } v\in V\\
  &\iff \ip{Tv}{w}=0 \text{ for all } v\in V\\
  &\iff w\in(\range T)^\perp.
\end{align*}
Thus $\nul T^*=(\range T)^\perp$, proving (a).

If we take the orthogonal complement of both sides of (a), we get (d), where we
have used \ref{ax:6.52}. Replacing $T$ with $T^*$ in (a) gives (c), where we
have used \ref{ax:7.5}(c). Finally, replacing $T$ with $T^*$ in (d) gives (b).
\end{proof}

As we will soon see, the next definition is intimately connected to the matrix
of the adjoint of a linear map.

\begin{definition}[Conjugate transpose, $A^*$]\label{ax:7.7}
The \emph{conjugate transpose} of an $m$-by-$n$ matrix $A$ is the $n$-by-$m$
matrix $A^*$ obtained by interchanging the rows and columns and then taking the
complex conjugate of each entry. In other words, if $j\in\{1,\dots,n\}$ and
$k\in\{1,\dots,m\}$, then
\[ (A^*)_{j,k} = \overline{A_{k,j}}. \]
\end{definition}

\begin{example}[Conjugate transpose of a $2$-by-$3$ matrix]\label{ax:7.8}
The conjugate transpose of the $2$-by-$3$ matrix
$\begin{pmatrix} 2 & 3+4i & 7\\ 6 & 5 & 8i \end{pmatrix}$
is the $3$-by-$2$ matrix
\[ \begin{pmatrix} 2 & 6\\ 3-4i & 5\\ 7 & -8i \end{pmatrix}. \]
\end{example}

% the first note sits beside Example 7.8 (a note cannot go inside the box)
\marginnote[-24mm]{If a matrix $A$ has only real entries, then $A^*=A^{\mathrm{t}}$,
where $A^{\mathrm{t}}$ denotes the transpose of $A$ (the matrix obtained by
interchanging the rows and the columns).}%
\marginnote{The adjoint of a linear map does not depend on a choice of basis.
Thus we frequently emphasize adjoints of linear maps instead of transposes or
conjugate transposes of matrices.}%
The next result shows how to compute the matrix of $T^*$ from the matrix of
$T$. \textbf{Caution:} With respect to nonorthonormal bases, the matrix of $T^*$
does not necessarily equal the conjugate transpose of the matrix of $T$.

\begin{result}[Matrix of $T^*$ equals conjugate transpose of matrix of $T$]%
\label{ax:7.9}
Let $T\in\Lin(V,W)$. Suppose $e_1,\dots,e_n$ is an orthonormal basis of $V$ and
$f_1,\dots,f_m$ is an orthonormal basis of $W$. Then
$\Mat\bigl(T^*,(f_1,\dots,f_m),(e_1,\dots,e_n)\bigr)$ is the conjugate transpose
of $\Mat\bigl(T,(e_1,\dots,e_n),(f_1,\dots,f_m)\bigr)$. In other words,
\[ \Mat(T^*) = \bigl(\Mat(T)\bigr)^*. \]
\end{result}

\begin{proof}
In this proof, we will write $\Mat(T)$ and $\Mat(T^*)$ instead of the longer
expressions $\Mat\bigl(T,(e_1,\dots,e_n),(f_1,\dots,f_m)\bigr)$ and
$\Mat\bigl(T^*,(f_1,\dots,f_m),(e_1,\dots,e_n)\bigr)$.

Recall that we obtain the $k^{\text{th}}$ column of $\Mat(T)$ by writing $Te_k$
as a linear combination of the $f_j$'s; the scalars used in this linear
combination then become the $k^{\text{th}}$ column of $\Mat(T)$. Because
$f_1,\dots,f_m$ is an orthonormal basis of $W$, we know how to write $Te_k$ as a
linear combination of the $f_j$'s [see \ref{ax:6.30}(a)]:
\[ Te_k = \ip{Te_k}{f_1}f_1+\cdots+\ip{Te_k}{f_m}f_m. \]
Thus
\[ \text{the entry in row $j$, column $k$, of $\Mat(T)$ is $\ip{Te_k}{f_j}$.} \]

In the statement above, replace $T$ with $T^*$ and interchange $e_1,\allowbreak\dots,e_n$
and $f_1,\allowbreak\dots,f_m$. This shows that the entry in row $j$, column $k$, of
$\Mat(T^*)$ is $\ip{T^*f_k}{e_j}$, which equals $\ip{f_k}{Te_j}$, which equals
$\overline{\ip{Te_j}{f_k}}$, which equals the complex conjugate of the entry in
row $k$, column $j$, of $\Mat(T)$. Thus $\Mat(T^*)=\bigl(\Mat(T)\bigr)^*$.
\end{proof}

The Riesz representation theorem as stated in \ref{ax:6.58} provides an
identification of $V$ with its dual space $V'$ defined in \ref{ax:3.110}. Under
this identification, the orthogonal complement $U^\perp$ of a subset
$U\subseteq V$ corresponds to the annihilator $U^0$ of $U$. If $U$ is a subspace
of $V$, then the formulas for the dimensions of $U^\perp$ and $U^0$ become
identical under this identification---see \ref{ax:3.125} and \ref{ax:6.51}.

\marginnote{Because orthogonal complements and adjoints are easier to deal with
than annihilators and dual maps, there is no need to work with annihilators and
dual maps in the context of inner product spaces.}%
Suppose $T\colon V\to W$ is a linear map. Under the identification of $V$ with
$V'$ and the identification of $W$ with $W'$, the adjoint map
$T^*\colon W\to V$ corresponds to the dual map $T'\colon W'\to V'$ defined in
\ref{ax:3.118}, as Exercise~32 asks you to verify. Under this identification,
the formulas for $\nul T^*$ and $\range T^*$ [\ref{ax:7.6}(a) and (b)] then
become identical to the formulas for $\nul T'$ and $\range T'$
[\ref{ax:3.128}(a) and \ref{ax:3.130}(b)]. Furthermore, the theorem about the
matrix of $T^*$ (\ref{ax:7.9}) is analogous to the theorem about the matrix of
$T'$ (\ref{ax:3.132}).

\subsection{Self-Adjoint Operators}

Now we switch our attention to operators on inner product spaces. Instead of
considering linear maps from $V$ to $W$, we will focus on linear maps from $V$
to $V$; recall that such linear maps are called operators.

\begin{definition}[Self-adjoint]\label{ax:7.10}
An operator $T\in\Lin(V)$ is called \emph{self-adjoint} if $T=T^*$.
\end{definition}

If $T\in\Lin(V)$ and $e_1,\dots,e_n$ is an orthonormal basis of $V$, then $T$
is self-adjoint if and only if
$\Mat\bigl(T,(e_1,\dots,e_n)\bigr)=\Mat\bigl(T,(e_1,\dots,e_n)\bigr)^*$, as
follows from \ref{ax:7.9}.

\begin{example}[Determining whether $T$ is self-adjoint from its matrix]%
\label{ax:7.11}
Suppose $c\in\F$ and $T$ is the operator on $\F^2$ whose matrix (with respect
to the standard basis) is
\[ \Mat(T) = \begin{pmatrix} 2 & c\\ 3 & 7 \end{pmatrix}. \]
The matrix of $T^*$ (with respect to the standard basis) is
\[ \Mat(T^*) = \begin{pmatrix} 2 & 3\\ \overline{c} & 7 \end{pmatrix}. \]
Thus $\Mat(T)=\Mat(T^*)$ if and only if $c=3$. Hence the operator $T$ is
self-adjoint if and only if $c=3$.
\end{example}

A good analogy to keep in mind is that the adjoint on $\Lin(V)$ plays a role
similar to that of the complex conjugate on $\C$. A complex number $z$ is real
if and only if $z=\overline{z}$; thus a self-adjoint operator ($T=T^*$) is
analogous to a real number.

\marginnote{An operator $T\in\Lin(V)$ is self-adjoint if and only if
\[ \ip{Tv}{w} = \ip{v}{Tw} \]
for all $v,w\in V$.}%
We will see that the analogy discussed above is reflected in some important
properties of self-adjoint operators, beginning with eigenvalues in the next
result.

If $\F=\R$, then by definition every eigenvalue is real, so the next result is
interesting only when $\F=\C$.

\begin{result}[Eigenvalues of self-adjoint operators]\label{ax:7.12}
Every eigenvalue of a self-adjoint operator is real.
\end{result}

\begin{proof}
Suppose $T$ is a self-adjoint operator on $V$. Let $\lambda$ be an eigenvalue of
$T$, and let $v$ be a nonzero vector in $V$ such that $Tv=\lambda v$. Then
\[ \lambda\norm{v}^2 = \ip{\lambda v}{v} = \ip{Tv}{v} = \ip{v}{Tv}
   = \ip{v}{\lambda v} = \overline{\lambda}\norm{v}^2. \]
Thus $\lambda=\overline{\lambda}$, which means that $\lambda$ is real, as
desired.
\end{proof}

The next result is false for real inner product spaces. As an example, consider
the operator $T\in\Lin(\R^2)$ that is a counterclockwise rotation of $90^\circ$
around the origin; thus $T(x,y)=(-y,x)$. Notice that $Tv$ is orthogonal to $v$
for every $v\in\R^2$, even though $T\ne0$.

\begin{result}[$Tv$ is orthogonal to $v$ for all $v$ $\iff$ $T=0$ (assuming
$\F=\C$)]\label{ax:7.13}
Suppose $V$ is a complex inner product space and $T\in\Lin(V)$. Then
\[ \ip{Tv}{v}=0 \text{ for every } v\in V \iff T=0. \]
\end{result}

\begin{proof}
If $u,w\in V$, then
\begin{align*}
\ip{Tu}{w} &= \frac{\bigl\langle T(u+w),u+w\bigr\rangle
                    -\bigl\langle T(u-w),u-w\bigr\rangle}{4}\\
  &\qquad+\frac{\bigl\langle T(u+iw),u+iw\bigr\rangle
                -\bigl\langle T(u-iw),u-iw\bigr\rangle}{4}\,i,
\end{align*}
as can be verified by computing the right side. Note that each term on the
right side is of the form $\ip{Tv}{v}$ for appropriate $v\in V$.

Now suppose $\ip{Tv}{v}=0$ for every $v\in V$. Then the equation above implies
that $\ip{Tu}{w}=0$ for all $u,w\in V$, which then implies that $Tu=0$ for every
$u\in V$ (take $w=Tu$). Hence $T=0$, as desired.
\end{proof}

\marginnote{The next result provides another good example of how self-adjoint
operators behave like real numbers.}%
The next result is false for real inner product spaces, as shown by considering
any operator on a real inner product space that is not self-adjoint.

\begin{result}[$\ip{Tv}{v}$ is real for all $v$ $\iff$ $T$ is self-adjoint
(assuming $\F=\C$)]\label{ax:7.14}
Suppose $V$ is a complex inner product space and $T\in\Lin(V)$. Then
\[ T \text{ is self-adjoint} \iff \ip{Tv}{v}\in\R \text{ for every } v\in V. \]
\end{result}

\begin{proof}
If $v\in V$, then\axeqstep{ax:7.15}
\[ \ip{T^*v}{v} = \overline{\ip{v}{T^*v}} = \overline{\ip{Tv}{v}}.
   \tag{\thetheorem} \]
Now
\begin{align*}
T \text{ is self-adjoint} &\iff T-T^*=0\\
  &\iff \bigl\langle (T-T^*)v,v\bigr\rangle=0 \text{ for every } v\in V\\
  &\iff \ip{Tv}{v}-\overline{\ip{Tv}{v}}=0 \text{ for every } v\in V\\
  &\iff \ip{Tv}{v}\in\R \text{ for every } v\in V,
\end{align*}
where the second equivalence follows from \ref{ax:7.13} as applied to $T-T^*$
and the third equivalence follows from \ref{ax:7.15}.
\end{proof}

On a real inner product space $V$, a nonzero operator $T$ might satisfy
$\ip{Tv}{v}=0$ for all $v\in V$. However, the next result shows that this
cannot happen for a self-adjoint operator.

\begin{result}[$T$ self-adjoint and $\ip{Tv}{v}=0$ for all $v$ $\iff$
$T=0$]\label{ax:7.16}
Suppose $T$ is a self-adjoint operator on $V$. Then
\[ \ip{Tv}{v}=0 \text{ for every } v\in V \iff T=0. \]
\end{result}

\begin{proof}
We have already proved this (without the hypothesis that $T$ is self-adjoint)
when $V$ is a complex inner product space (see \ref{ax:7.13}). Thus we can
assume that $V$ is a real inner product space. If $u,w\in V$, then%
\axeqstep{ax:7.17}
\[ \ip{Tu}{w} = \frac{\bigl\langle T(u+w),u+w\bigr\rangle
                      -\bigl\langle T(u-w),u-w\bigr\rangle}{4},
   \tag{\thetheorem} \]
as can be proved by computing the right side using the equation
\[ \ip{Tw}{u} = \ip{w}{Tu} = \ip{Tu}{w}, \]
where the first equality holds because $T$ is self-adjoint and the second
equality holds because we are working in a real inner product space.

Now suppose $\ip{Tv}{v}=0$ for every $v\in V$. Because each term on the right
side of \ref{ax:7.17} is of the form $\ip{Tv}{v}$ for appropriate $v$, this
implies that $\ip{Tu}{w}=0$ for all $u,w\in V$. This implies that $Tu=0$ for
every $u\in V$ (take $w=Tu$). Hence $T=0$, as desired.
\end{proof}

\subsection{Normal Operators}

\begin{definition}[Normal]\label{ax:7.18}\leavevmode
\begin{itemize}
\item An operator on an inner product space is called \emph{normal} if it
  commutes with its adjoint.
\item In other words, $T\in\Lin(V)$ is normal if $TT^*=T^*T$.
\end{itemize}
\end{definition}

Every self-adjoint operator is normal, because if $T$ is self-adjoint then
$T^*=T$ and hence $T$ commutes with $T^*$.

\begin{example}[An operator that is normal but not self-adjoint]\label{ax:7.19}
Let $T$ be the operator on $\F^2$ whose matrix (with respect to the standard
basis) is
\[ \begin{pmatrix} 2 & -3\\ 3 & 2 \end{pmatrix}. \]
Thus $T(w,z)=(2w-3z,3w+2z)$.

This operator $T$ is not self-adjoint because the entry in row~2, column~1
(which equals $3$) does not equal the complex conjugate of the entry in row~1,
column~2 (which equals $-3$).

The matrix of $TT^*$ equals
\[ \begin{pmatrix} 2 & -3\\ 3 & 2 \end{pmatrix}
   \begin{pmatrix} 2 & 3\\ -3 & 2 \end{pmatrix},
   \quad\text{which equals}\quad
   \begin{pmatrix} 13 & 0\\ 0 & 13 \end{pmatrix}. \]
Similarly, the matrix of $T^*T$ equals
\[ \begin{pmatrix} 2 & 3\\ -3 & 2 \end{pmatrix}
   \begin{pmatrix} 2 & -3\\ 3 & 2 \end{pmatrix},
   \quad\text{which equals}\quad
   \begin{pmatrix} 13 & 0\\ 0 & 13 \end{pmatrix}. \]
Because $TT^*$ and $T^*T$ have the same matrix, we see that $TT^*=T^*T$. Thus
$T$ is normal.
\end{example}

In the next section we will see why normal operators are worthy of special
attention. The next result provides a useful characterization of normal
operators.

\begin{result}[$T$ is normal if and only if $Tv$ and $T^*v$ have the same
norm]\label{ax:7.20}
Suppose $T\in\Lin(V)$. Then
\[ T \text{ is normal} \iff \norm{Tv}=\norm{T^*v} \text{ for every } v\in V. \]
\end{result}

\begin{proof}
We have
\begin{align*}
T \text{ is normal} &\iff T^*T-TT^*=0\\
  &\iff \bigl\langle (T^*T-TT^*)v,v\bigr\rangle=0 \text{ for every } v\in V\\
  &\iff \ip{T^*Tv}{v}=\ip{TT^*v}{v} \text{ for every } v\in V\\
  &\iff \ip{Tv}{Tv}=\ip{T^*v}{T^*v} \text{ for every } v\in V\\
  &\iff \norm{Tv}^2=\norm{T^*v}^2 \text{ for every } v\in V\\
  &\iff \norm{Tv}=\norm{T^*v} \text{ for every } v\in V,
\end{align*}
where we used \ref{ax:7.16} to establish the second equivalence (note that the
operator $T^*T-TT^*$ is self-adjoint).
\end{proof}

The next result presents several consequences of the result above. Compare (e)
of the next result to Exercise~3. That exercise states that the eigenvalues of
the adjoint of each operator are equal (as a set) to the complex conjugates of
the eigenvalues of the operator. The exercise says nothing about eigenvectors,
because an operator and its adjoint may have different eigenvectors. However,
(e) of the next result implies that a normal operator and its adjoint have the
same eigenvectors.

\begin{result}[Range, null space, and eigenvectors of a normal operator]%
\label{ax:7.21}
Suppose $T\in\Lin(V)$ is normal. Then
\begin{parts}
\item $\nul T=\nul T^*$;
\item $\range T=\range T^*$;
\item $V=\nul T\oplus\range T$;
\item $T-\lambda I$ is normal for every $\lambda\in\F$;
\item if $v\in V$ and $\lambda\in\F$, then $Tv=\lambda v$ if and only if
  $T^*v=\overline{\lambda}v$.
\end{parts}
\end{result}

\begin{proof}\leavevmode
\begin{parts}
\item Suppose $v\in V$. Then
  \[ v\in\nul T \iff \norm{Tv}=0 \iff \norm{T^*v}=0 \iff v\in\nul T^*, \]
  where the middle equivalence above follows from \ref{ax:7.20}. Thus
  $\nul T=\nul T^*$.
\item We have
  \[ \range T = (\nul T^*)^\perp = (\nul T)^\perp = \range T^*, \]
  where the first equality comes from \ref{ax:7.6}(d), the second equality
  comes from (a) in this result, and the third equality comes from
  \ref{ax:7.6}(b).
\item We have
  \[ V = (\nul T)\oplus(\nul T)^\perp = \nul T\oplus\range T^*
       = \nul T\oplus\range T, \]
  where the first equality comes from \ref{ax:6.49}, the second equality comes
  from \ref{ax:7.6}(b), and the third equality comes from (b) in this result.
\item Suppose $\lambda\in\F$. Then
  \begin{align*}
  (T-\lambda I)(T-\lambda I)^* &= (T-\lambda I)(T^*-\overline{\lambda}I)\\
    &= TT^*-\overline{\lambda}T-\lambda T^*+\abs{\lambda}^2I\\
    &= T^*T-\overline{\lambda}T-\lambda T^*+\abs{\lambda}^2I\\
    &= (T^*-\overline{\lambda}I)(T-\lambda I)\\
    &= (T-\lambda I)^*(T-\lambda I).
  \end{align*}
  Thus $T-\lambda I$ commutes with its adjoint. Hence $T-\lambda I$ is normal.
\item Suppose $v\in V$ and $\lambda\in\F$. Then (d) and \ref{ax:7.20} imply that
  \[ \norm{(T-\lambda I)v} = \norm{(T-\lambda I)^*v}
     = \norm{(T^*-\overline{\lambda}I)v}. \]
  Thus $\norm{(T-\lambda I)v}=0$ if and only if
  $\norm{(T^*-\overline{\lambda}I)v}=0$. Hence $Tv=\lambda v$ if and only if
  $T^*v=\overline{\lambda}v$. \qedhere
\end{parts}
\end{proof}

Because every self-adjoint operator is normal, the next result applies in
particular to self-adjoint operators.

\begin{result}[Orthogonal eigenvectors for normal operators]\label{ax:7.22}
Suppose $T\in\Lin(V)$ is normal. Then eigenvectors of $T$ corresponding to
distinct eigenvalues are orthogonal.
\end{result}

\begin{proof}
Suppose $\alpha,\beta$ are distinct eigenvalues of $T$, with corresponding
eigenvectors $u,v$. Thus $Tu=\alpha u$ and $Tv=\beta v$. From \ref{ax:7.21}(e)
we have $T^*v=\overline{\beta}v$. Thus
\begin{align*}
(\alpha-\beta)\ip{u}{v} &= \ip{\alpha u}{v}-\ip{u}{\overline{\beta}v}\\
                        &= \ip{Tu}{v}-\ip{u}{T^*v}\\
                        &= 0.
\end{align*}
Because $\alpha\ne\beta$, the equation above implies that $\ip{u}{v}=0$. Thus
$u$ and $v$ are orthogonal, as desired.
\end{proof}

As stated here, the next result makes sense only when $\F=\C$. However, see
Exercise~12 for a version that makes sense when $\F=\C$ and when $\F=\R$.

Suppose $\F=\C$ and $T\in\Lin(V)$. Under the analogy between $\Lin(V)$ and
$\C$, with the adjoint on $\Lin(V)$ playing a similar role to that of the
complex conjugate on $\C$, the operators $A$ and $B$ as defined by \ref{ax:7.24}
correspond to the real and imaginary parts of $T$. Thus the informal title of
the result below should make sense.

\begin{result}[$T$ is normal $\iff$ the real and imaginary parts of $T$
commute]\label{ax:7.23}
Suppose $\F=\C$ and $T\in\Lin(V)$. Then $T$ is normal if and only if there
exist commuting self-adjoint operators $A$ and $B$ such that $T=A+iB$.
\end{result}

\begin{proof}
First suppose $T$ is normal. Let\axeqstep{ax:7.24}
\[ A = \frac{T+T^*}{2} \quad\text{and}\quad B = \frac{T-T^*}{2i}.
   \tag{\thetheorem} \]
Then $A$ and $B$ are self-adjoint and $T=A+iB$. A quick computation shows
that\axeqstep{ax:7.25}
\[ AB-BA = \frac{T^*T-TT^*}{2i}. \tag{\thetheorem} \]
Because $T$ is normal, the right side of the equation above equals $0$. Thus
the operators $A$ and $B$ commute, as desired.

To prove the implication in the other direction, now suppose there exist
commuting self-adjoint operators $A$ and $B$ such that $T=A+iB$. Then
$T^*=A-iB$. Adding the last two equations and then dividing by $2$ produces the
equation for $A$ in \ref{ax:7.24}. Subtracting the last two equations and then
dividing by $2i$ produces the equation for $B$ in \ref{ax:7.24}. Now
\ref{ax:7.24} implies \ref{ax:7.25}. Because $B$ and $A$ commute,
\ref{ax:7.25} implies that $T$ is normal, as desired.
\end{proof}

% ------------------------------------------------------------------ exercises
\exercisesheading{7A}

\begin{exercise}
Suppose $n$ is a positive integer. Define $T\in\Lin(\F^n)$ by
\[ T(z_1,\dots,z_n) = (0,z_1,\dots,z_{n-1}). \]
Find a formula for $T^*(z_1,\dots,z_n)$.
\end{exercise}

\begin{exercise}
Suppose $T\in\Lin(V,W)$. Prove that
\[ T=0 \iff T^*=0 \iff T^*T=0 \iff TT^*=0. \]
\end{exercise}

\begin{exercise}
Suppose $T\in\Lin(V)$ and $\lambda\in\F$. Prove that
\[ \lambda \text{ is an eigenvalue of } T \iff \overline{\lambda} \text{ is an
   eigenvalue of } T^*. \]
\end{exercise}

\begin{exercise}
Suppose $T\in\Lin(V)$ and $U$ is a subspace of $V$. Prove that
\[ U \text{ is invariant under } T \iff U^\perp \text{ is invariant under }
   T^*. \]
\end{exercise}

\begin{exercise}
Suppose $T\in\Lin(V,W)$. Suppose $e_1,\dots,e_n$ is an orthonormal basis of $V$
and $f_1,\dots,f_m$ is an orthonormal basis of $W$. Prove that
\[ \norm{Te_1}^2+\cdots+\norm{Te_n}^2 = \norm{T^*f_1}^2+\cdots+\norm{T^*f_m}^2. \]
\exercisenote{The numbers $\norm{Te_1}^2,\dots,\norm{Te_n}^2$ in the equation
above depend on the orthonormal basis $e_1,\dots,e_n$, but the right side of the
equation does not depend on $e_1,\dots,e_n$. Thus the equation above shows that
the sum on the left side does not depend on which orthonormal basis
$e_1,\dots,e_n$ is used.}
\end{exercise}

\begin{exercise}
Suppose $T\in\Lin(V,W)$. Prove that
\begin{parts}
\item $T$ is injective $\iff$ $T^*$ is surjective;
\item $T$ is surjective $\iff$ $T^*$ is injective.
\end{parts}
\end{exercise}

\begin{exercise}
Prove that if $T\in\Lin(V,W)$, then
\begin{parts}
\item $\dim\nul T^*=\dim\nul T+\dim W-\dim V$;
\item $\dim\range T^*=\dim\range T$.
\end{parts}
\end{exercise}

\begin{exercise}
Suppose $A$ is an $m$-by-$n$ matrix with entries in $\F$. Use (b) in
Exercise~7 to prove that the row rank of $A$ equals the column rank of $A$.
\exercisenote{This exercise asks for yet another alternative proof of a result
that was previously proved in \ref{ax:3.57} and \ref{ax:3.133}.}
\end{exercise}

\begin{exercise}
Prove that the product of two self-adjoint operators on $V$ is self-adjoint if
and only if the two operators commute.
\end{exercise}

\begin{exercise}
Suppose $\F=\C$ and $T\in\Lin(V)$. Prove that $T$ is self-adjoint if and only
if
\[ \ip{Tv}{v} = \bigl\langle T^*v,v\bigr\rangle \]
for all $v\in V$.
\end{exercise}

\begin{exercise}
Define an operator $S\colon\F^2\to\F^2$ by $S(w,z)=(-z,w)$.
\begin{parts}
\item Find a formula for $S^*$.
\item Show that $S$ is normal but not self-adjoint.
\item Find all eigenvalues of $S$.
\end{parts}
\exercisenote{If $\F=\R$, then $S$ is the operator on $\R^2$ of
counterclockwise rotation by $90^\circ$.}
\end{exercise}

\begin{exercise}
An operator $B\in\Lin(V)$ is called \emph{skew} if
\[ B^* = -B. \]
Suppose that $T\in\Lin(V)$. Prove that $T$ is normal if and only if there exist
commuting operators $A$ and $B$ such that $A$ is self-adjoint, $B$ is a skew
operator, and $T=A+B$.
\end{exercise}

\begin{exercise}
Suppose $\F=\R$. Define $\mathcal{A}\in\Lin\bigl(\Lin(V)\bigr)$ by
$\mathcal{A}T=T^*$ for all $T\in\Lin(V)$.
\begin{parts}
\item Find all eigenvalues of $\mathcal{A}$.
\item Find the minimal polynomial of $\mathcal{A}$.
\end{parts}
\end{exercise}

\begin{exercise}
Define an inner product on $\Poly_2(\R)$ by $\ip{p}{q}=\int_0^1pq$. Define an
operator $T\in\Lin\bigl(\Poly_2(\R)\bigr)$ by
\[ T(ax^2+bx+c) = bx. \]
\begin{parts}
\item Show that with this inner product, the operator $T$ is not self-adjoint.
\item The matrix of $T$ with respect to the basis $1,x,x^2$ is
  \[ \begin{pmatrix} 0 & 0 & 0\\ 0 & 1 & 0\\ 0 & 0 & 0 \end{pmatrix}. \]
  This matrix equals its conjugate transpose, even though $T$ is not
  self-adjoint. Explain why this is not a contradiction.
\end{parts}
\end{exercise}

\begin{exercise}
Suppose $T\in\Lin(V)$ is invertible. Prove that
\begin{parts}
\item $T$ is self-adjoint $\iff$ $T^{-1}$ is self-adjoint;
\item $T$ is normal $\iff$ $T^{-1}$ is normal.
\end{parts}
\end{exercise}

\begin{exercise}
Suppose $\F=\R$.
\begin{parts}
\item Show that the set of self-adjoint operators on $V$ is a subspace of
  $\Lin(V)$.
\item What is the dimension of the subspace of $\Lin(V)$ in (a) [in terms of
  $\dim V$]?
\end{parts}
\end{exercise}

\begin{exercise}
Suppose $\F=\C$. Show that the set of self-adjoint operators on $V$ is not a
subspace of $\Lin(V)$.
\end{exercise}

\begin{exercise}
Suppose $\dim V\ge2$. Show that the set of normal operators on $V$ is not a
subspace of $\Lin(V)$.
\end{exercise}

\begin{exercise}
Suppose $T\in\Lin(V)$ and $\norm{T^*v}\le\norm{Tv}$ for every $v\in V$. Prove
that $T$ is normal.
\exercisenote{This exercise fails on infinite-dimensional inner product spaces,
leading to what are called hyponormal operators, which have a well-developed
theory.}
\end{exercise}

\begin{exercise}
Suppose $P\in\Lin(V)$ is such that $P^2=P$. Prove that the following are
equivalent.
\begin{parts}
\item $P$ is self-adjoint.
\item $P$ is normal.
\item There is a subspace $U$ of $V$ such that $P=P_U$.
\end{parts}
\end{exercise}

\begin{exercise}
Suppose $D\colon\Poly_8(\R)\to\Poly_8(\R)$ is the differentiation operator
defined by $Dp=p'$. Prove that there does not exist an inner product on
$\Poly_8(\R)$ that makes $D$ a normal operator.
\end{exercise}

\begin{exercise}
Give an example of an operator $T\in\Lin(\R^3)$ such that $T$ is normal but not
self-adjoint.
\end{exercise}

\begin{exercise}
Suppose $T$ is a normal operator on $V$. Suppose also that $v,w\in V$ satisfy
the equations
\[ \norm{v}=\norm{w}=2, \qquad Tv=3v, \qquad Tw=4w. \]
Show that $\norm{T(v+w)}=10$.
\end{exercise}

\begin{exercise}
Suppose $T\in\Lin(V)$ and
\[ a_0+a_1z+a_2z^2+\cdots+a_{m-1}z^{m-1}+z^m \]
is the minimal polynomial of $T$. Prove that the minimal polynomial of $T^*$ is
\[ \overline{a_0}+\overline{a_1}z+\overline{a_2}z^2+\cdots
   +\overline{a_{m-1}}z^{m-1}+z^m. \]
\exercisenote{This exercise shows that the minimal polynomial of $T^*$ equals
the minimal polynomial of $T$ if $\F=\R$.}
\end{exercise}

\begin{exercise}
Suppose $T\in\Lin(V)$. Prove that $T$ is diagonalizable if and only if $T^*$ is
diagonalizable.
\end{exercise}

\begin{exercise}
Fix $u,x\in V$. Define $T\in\Lin(V)$ by $Tv=\ip{v}{u}x$ for every $v\in V$.
\begin{parts}
\item Prove that if $V$ is a real vector space, then $T$ is self-adjoint if and
  only if the list $u,x$ is linearly dependent.
\item Prove that $T$ is normal if and only if the list $u,x$ is linearly
  dependent.
\end{parts}
\end{exercise}

\begin{exercise}
Suppose $T\in\Lin(V)$ is normal. Prove that
\[ \nul T^k=\nul T \quad\text{and}\quad \range T^k=\range T \]
for every positive integer $k$.
\end{exercise}

\begin{exercise}
Suppose $T\in\Lin(V)$ is normal. Prove that if $\lambda\in\F$, then the minimal
polynomial of $T$ is not a polynomial multiple of $(x-\lambda)^2$.
\end{exercise}

\begin{exercise}
Prove or give a counterexample: If $T\in\Lin(V)$ and there is an orthonormal
basis $e_1,\dots,e_n$ of $V$ such that $\norm{Te_k}=\norm{T^*e_k}$ for each
$k=1,\dots,n$, then $T$ is normal.
\end{exercise}

\begin{exercise}
Suppose that $T\in\Lin(\F^3)$ is normal and $T(1,1,1)=(2,2,2)$. Suppose
$(z_1,z_2,z_3)\in\nul T$. Prove that $z_1+z_2+z_3=0$.
\end{exercise}

\begin{exercise}
Fix a positive integer $n$. In the inner product space of continuous
real-valued functions on $[-\pi,\pi]$ with inner product
$\ip{f}{g}=\int_{-\pi}^{\pi}fg$, let
\[ V = \spn(1,\cos x,\cos2x,\dots,\cos nx,\sin x,\sin2x,\dots,\sin nx). \]
\begin{parts}
\item Define $D\in\Lin(V)$ by $Df=f'$. Show that $D^*=-D$. Conclude that $D$ is
  normal but not self-adjoint.
\item Define $T\in\Lin(V)$ by $Tf=f''$. Show that $T$ is self-adjoint.
\end{parts}
\end{exercise}

\begin{exercise}
Suppose $T\colon V\to W$ is a linear map. Show that under the standard
identification of $V$ with $V'$ (see \ref{ax:6.58}) and the corresponding
identification of $W$ with $W'$, the adjoint map $T^*\colon W\to V$ corresponds
to the dual map $T'\colon W'\to V'$. More precisely, show that
\[ T'(\varphi_w) = \varphi_{T^*w} \]
for all $w\in W$, where $\varphi_w$ and $\varphi_{T^*w}$ are defined as in
\ref{ax:6.58}.
\end{exercise}
